\documentclass[11pt]{article}

% Shared preamble for the SKJS soundness corpus (main paper + deep dives).
% Loaded via \input in each document's preamble. Documents may define
% additional packages, theorem environments, and macros AFTER this input,
% but must not redefine anything here.
\usepackage{amsmath}
\usepackage{amssymb}
\usepackage{amsthm}
\usepackage{mathtools}
\usepackage{stmaryrd}
\usepackage[margin=1in]{geometry}
\usepackage{tcolorbox}
\usepackage{booktabs}
\usepackage{array}
\usepackage{longtable}
\usepackage[hyphens]{url}
\urlstyle{tt}

% Long typewriter tokens (paths, identifiers) defeat hyphenation;
% allow inter-word stretch instead of overfull lines.
\emergencystretch=2em

\newtheorem{theorem}{Theorem}
\newtheorem{lemma}[theorem]{Lemma}
\newtheorem{definition}[theorem]{Definition}
\newtheorem{corollary}[theorem]{Corollary}
\newtheorem{axiom}[theorem]{Axiom}
\newtheorem{obligation}[theorem]{Obligation}

% Summary box for key definitions
\newtcolorbox{keydef}{
  colback=gray!5,
  colframe=gray!50,
  boxrule=0.5pt,
  arc=2pt,
  left=4pt,right=4pt,top=2pt,bottom=2pt
}

% Core semantic notation shared by every document
\newcommand{\sem}[1]{\llbracket #1 \rrbracket}
\newcommand{\Val}{\mathsf{Val}}

% Implementation anchor: file / class / function at the pinned revision
\newcommand{\impl}[1]{\par\noindent{\footnotesize\textsf{impl:} \texttt{#1}}}

% Single source for the modeled revision. Every document's \date uses this;
% the ledger and tooling read the same values.
\newcommand{\skjspinhash}{9108f339}
\newcommand{\skjspindate}{2026-07-02}
\newcommand{\skjspin}{Models revision \texttt{\skjspinhash} of \texttt{main} (\skjspindate)}



\newtheorem{fact}[theorem]{Fact}
\newtheorem{finding}[theorem]{Finding}

% Notation copied from the main paper (skjs-type-soundness.tex)
\newcommand{\Payload}{\mathsf{Payload}}
\newcommand{\String}{\mathsf{String}}
\newcommand{\SymVal}{\mathsf{SymVal}}
\newcommand{\tagfn}{\mathsf{tag}}
\newcommand{\rtag}{\mathsf{rtag}}


% Local notation
\newcommand{\Lit}[1]{\mathsf{Lit}_{\texttt{#1}}}
\newcommand{\Canon}[1]{\mathsf{Canon}_{\texttt{#1}}}
\newcommand{\canon}{\mathsf{canon}}
\newcommand{\render}{\mathsf{render}}
\newcommand{\Fdb}{\mathbb{F}}
\newcommand{\zeq}{\doteq}   % IEEE equality with -0 identified with 0
\newcommand{\truthy}{\mathsf{truthy}}

\title{Literal Canonicalization in SKJS:\\
Discharging Assumption \textbf{L-LIT} (Deep Dive G)}
\author{}
\date{\skjspin.\\
Companion to \emph{Refinement and Inference Growth in SKJS} (the main paper).}

\begin{document}

\maketitle

\begin{abstract}
The main soundness paper treats literal types abstractly through a
denotation $\delta_k$ from stored literal spellings to runtime payloads, and
records three claims about it as assumption \textbf{L-LIT}: (i) $\delta$ is
a function, (ii) distinct (normalized) spellings denote
$\approx$-inequivalent values, and (iii) the implementation's lexeme-level
falsiness test agrees with the semantic definition of truthiness. At the
pinned revision, parts (ii) and (iii) are \textsc{violated}: numeric and
bigint literal types store the raw source lexeme, so the value $0$ has
unboundedly many spellings, of which the stopgap
\texttt{isZeroSpelling} recognizes exactly~14 (issue \#655). This deep dive
is the fix specification and the mathematics that supports it. We define
the denotation $\delta_k$ precisely for each literal kind, define a
canonical form per kind (parse to the denoted value, re-render canonically,
with $-0$ identified with $0$ to match the main paper's $\approx$-closure
semantics), and prove that over canonical forms \textbf{L-LIT}(i)--(iii)
all hold, with the exact finite falsy canonical set per kind. We then
specify where normalization should live, what it does to user-visible
pretty-printing, how it interacts with \texttt{PrimLiterals} set
operations and numeric property names, what code becomes deletable, and
which pinned tests should accompany the change. Two adjacent gaps found
while attempting the proofs are reported as findings: undecoded string
escape sequences — a second, string-kind violation of \textbf{L-LIT}(ii),
witnessed end-to-end (Finding~\ref{find:escapes}) — and an
exponent-threshold over-acceptance in \texttt{isJSNumericString},
verified to have no soundness leak (hygiene only).
\end{abstract}

\tableofcontents

%==========================================================================
\section{Purpose}
\label{sec:purpose}

This paper serves the assumptions-and-obligations ledger of the main paper
(its Section ``Assumptions and Proof Obligations'') as follows.

\begin{itemize}
\item \textbf{L-LIT(i)} ($\delta$ is a function): currently marked
  \textsc{proved} (definitional). This paper \emph{strengthens} the row: we
  give $\delta_k$ concretely for every kind
  (Section~\ref{sec:denotations}), prove that no numeric literal denotes
  $\mathsf{NaN}$ (which the main paper's Definition of Literal Denotation
  assumes), and expose one caveat for kind \texttt{str} (escape decoding,
  Finding~\ref{find:escapes}).
\item \textbf{L-LIT(ii)} (injectivity modulo $\approx$): currently
  \textsc{violated} by \#655. This paper defines the canonical forms and
  \emph{proves} injectivity-modulo-$\approx$ over them
  (Theorem~\ref{thm:inj}). The row becomes \textsc{proved} once the fix of
  Section~\ref{sec:fix} lands; until then the \textsc{violated} marking
  stands.
\item \textbf{L-LIT(iii)} (falsiness-exactness): currently
  \textsc{violated} by \#655, with an observable unsoundness witness
  (main paper, worked example ``Known unsoundness: non-canonical zero
  spellings''). This paper proves falsiness-exactness over canonical
  forms (Theorem~\ref{thm:falsy}), giving the exact finite falsy canonical
  set per kind; same status trajectory as (ii).
\end{itemize}

Downstream, discharging (ii) and (iii) unconditionally restores: the
literal-filtering case of the main paper's Lemma ``Residues are Sound''
(the one soundness-relevant caveat in that document), and the exactness
clause of its Theorem ``Soundness of Constant Removal'' (precision only).
No theorem statement or proof in the main paper changes; the violations
are local to the leaf assumption, exactly as its summary asserts.

The fix itself is \emph{not} implemented at revision \texttt{9108f339};
Section~\ref{sec:fix} is its specification and Section~\ref{sec:status}
gives the precise ledger-row edits for both the interim and the post-fix
states.

%==========================================================================
\section{Recap of Definitions Used}
\label{sec:recap}

We reuse the main paper's runtime model verbatim and re-state only what the
proofs below manipulate; everything else is cited by name.

\begin{keydef}
From the main paper, Section ``Runtime Model'':
\begin{itemize}
\item Values are tagged payloads,
  $\Val = \{(k,p) \mid k \in \mathsf{Tags},\ p \in \Payload_k\}$, with
  $\Payload_{\mathsf{num}} = \Fdb$ (IEEE-754 doubles, including $-0$ and
  $\mathsf{NaN}$), $\Payload_{\mathsf{str}} = \String$,
  $\Payload_{\mathsf{bigint}} = \mathbb{Z}$,
  $\Payload_{\mathsf{bool}} = \{\mathsf{true},\mathsf{false}\}$,
  $\Payload_{\mathsf{sym}} = \SymVal$.
\item Strict equality $\approx$: same tag and equal payloads, except that
  on numbers $-0$ is identified with $0$ and $\mathsf{NaN}$ is equal to
  nothing, and the $\mathsf{undef}/\mathsf{void}$ tags are related
  cross-tag (irrelevant here: no literal type has a unit tag).
\item Truthiness (main paper, Definition ``Truthiness''):
  $(\mathsf{num}, n)$ is falsy iff $n \in \{0, -0, \mathsf{NaN}\}$;
  $(\mathsf{str}, s)$ falsy iff $s = \texttt{""}$;
  $(\mathsf{bigint}, z)$ falsy iff $z = 0$;
  $(\mathsf{bool}, b)$ falsy iff $b = \mathsf{false}$;
  symbols always truthy.
\item Literal types: $\mathsf{prim}(k, L)$ with
  $L \subseteq \Lit{k}$ finite nonempty, and the \emph{same-tag
  $\approx$-closure} semantics
  \[
  \sem{\mathsf{prim}(k, L)} =
  \{\, v \mid \rtag(v) = \tagfn(k),\
     \exists \ell \in L.\ v \approx (\tagfn(k), \delta_k(\ell)) \,\}
  \]
  so the literal type \texttt{0} contains both $0$ and $-0$, and
  $\Lit{num}$ is assumed to contain no $\mathsf{NaN}$ spelling.
\end{itemize}
\end{keydef}

For numbers we write $\zeq$ for the restriction of $\approx$ to
$\mathsf{num}$ payloads: $x \zeq y$ iff $x$ and $y$ are equal as IEEE
doubles with $-0$ identified with $0$; $\mathsf{NaN} \not\zeq
\mathsf{NaN}$. On $\Fdb \setminus \{\mathsf{NaN}\}$, $\zeq$ is an
equivalence whose only non-singleton class is $\{0, -0\}$.

\paragraph{Assumption L-LIT, as stated in the main paper's ledger.}
(i)~$\delta$ is a function --- definitional, used by the unconditional part
(R2) of constant-removal soundness; (ii)~\emph{injectivity modulo
$\approx$} after normalization --- distinct normalized lexemes denote
$\approx$-inequivalent values; needed only for \emph{exactness} claims;
(iii)~\emph{falsiness-exactness} --- the lexeme test \texttt{isFalsy}
agrees with the Definition of Truthiness on every stored literal; needed
by the residue lemma. Anchor: \texttt{PrimLiterals.sk};
\texttt{isZeroSpelling}. Status at \texttt{9108f339}: (i) \textsc{proved},
(ii)/(iii) \textsc{violated} by \#655.

%==========================================================================
\section{The Implementation Today, and the Gap}
\label{sec:today}

\subsection{Storage and comparison}

A literal is stored as a bare string with a position; identity, ordering,
hashing, and all set operations are by that string:
\texttt{PrimLiteral(pos, val)} projects to \texttt{val} and compares by
it; \texttt{PrimLiterals} is a string-keyed array map whose
\texttt{union}/\texttt{inter}/\texttt{diff}/\texttt{contains} are exact
set operations \emph{on lexemes}, and whose \texttt{xxhash} hashes only
the strings (positions excluded).
\impl{PrimLiterals.sk / PrimLiteral, PrimLiterals (union, inter, diff,
contains, xxhash)}

The kind (\texttt{num}, \texttt{bigint}, \texttt{str}, \texttt{bool},
\texttt{sym}) lives one level up, on
\texttt{TPrimValue(kind, literals)}; \texttt{PrimLiteral} itself is
kind-blind.
\impl{TType.sk / TPrimValue, TPrimValueKind}

\subsection{Construction sites of literal lexemes}
\label{sec:sites}

Every path that mints a numeric or bigint literal lexeme at
\texttt{9108f339}:

\begin{center}
\small
\renewcommand{\arraystretch}{1.2}
\begin{tabular}{@{}p{5.6cm}p{7.6cm}@{}}
\toprule
\textbf{Site} & \textbf{Lexeme minted} \\
\midrule
\texttt{ParserUtils.sk / getNumber, getBigInt} &
  raw token text, with an explicit \texttt{TODO: check and normalize
  somehow} \\
\texttt{ParserType.sk / parsePrimary} (lines 69--70, 232--233) &
  type annotations \texttt{TANumLiteral}/\texttt{TABigIntLiteral};
  negative forms prepend \texttt{"-"} to the raw lexeme \\
\texttt{TEnv.sk / normalize} (lines 1054, 1205) &
  converts those AST literals to \texttt{TPrimValue::single} unchanged \\
\texttt{TypeCheckExpr.sk / directLiteralSType} &
  expression literals (\texttt{Number}, \texttt{BigInt} tokens); unary
  \texttt{-} prepends \texttt{"-"} \\
\texttt{TypeCheckExpr.sk / asConstExpr} (\texttt{UOMinus} arm) &
  sign-flips by string prefix manipulation \\
\texttt{TEnv.sk} line 2162 (\texttt{keyof}) &
  string keys passing \texttt{isJSNumericString} become number literals
  with the same string \\
\texttt{SyntheticTuples.sk} line 127 &
  tuple arity via \texttt{Int.toString()} (canonical decimal) \\
\texttt{TType.sk / truthinessValueResidue} &
  the constants \texttt{"0n"}, \texttt{""} \\
\bottomrule
\end{tabular}
\end{center}

No site normalizes. Consequently the stored lexeme set for a source union
\texttt{0x1 | 1} has \emph{two} elements although the type is semantically
a singleton.

\subsection{The stopgap falsiness test}
\label{sec:stopgap}

Truthiness classification of a literal is by lexeme:
\texttt{TPrimValueKind::isFalsy} answers \texttt{literal == "false"}
(bool), \texttt{literal == ""} (str), always-false (sym), and for
num/bigint delegates to the stopgap \texttt{isZeroSpelling}
(TType.sk, ${\sim}$line 1242), which recognizes \emph{exactly 14}
spellings: signed decimal \texttt{0}/\texttt{-0} and the radix-prefixed
single zeros \texttt{0x0}/\texttt{0X0}/\texttt{0b0}/\texttt{0B0}/%
\texttt{0o0}/\texttt{0O0} with optional sign. The comment at
${\sim}$line 1235 marks this as a stopgap for issue \#655: floats
(\texttt{0.0}, \texttt{0e0}), digit separators (\texttt{0\_0}),
multi-digit radix zeros (\texttt{0x00}), and \texttt{00} are all
misclassified as truthy, and the same for their bigint forms. The num arm
additionally treats the lexeme \texttt{"NaN"} as falsy --- defensive only,
since no construction site of Section~\ref{sec:sites} can mint it
(Theorem~\ref{thm:nonan} makes this precise).
\impl{TType.sk / TPrimValueKind::isFalsy (${\sim}$1224),
isZeroSpelling (${\sim}$1242)}

Three consumers propagate the misclassification:
\texttt{SType::falsiness} counts falsy literals per member
(feeding the \texttt{safe}-mode condition guard in
\texttt{TypeCheckUtils.sk / checkCondition}, which requires
``boolean or a type with a single falsy value'');
\texttt{SType::truthinessResidue} filters literals by
\texttt{isFalsy} for \texttt{if}/\texttt{else} narrowing; and
\texttt{truthinessValueResidue} does the same for the value of
\texttt{\&\&}/\texttt{||}.
\impl{TType.sk / SType::falsiness, truthinessResidue,
truthinessValueResidue; TypeCheckUtils.sk / checkCondition}

\subsection{Observed witnesses}
\label{sec:witnesses}

The following behaviors were verified against the built checker at
\texttt{9108f339} (truthiness mode \texttt{safe}); this paper treats them
as established facts, and the main paper pins the first two in its worked
example ``Known unsoundness: non-canonical zero spellings (\#655)''.

\begin{enumerate}
\item \emph{Unsound acceptance.} For \texttt{x: 0.0 | ""}, both members
  are falsy at runtime, but \texttt{0.0} is classified truthy, so the
  union appears single-falsy and \texttt{if (x)} is accepted; the falsy
  residue drops \texttt{0.0}, and in the else-branch
\begin{verbatim}
function f(x: 0.0 | "") {
  if (x) {} else { const z: "" = x; }  // ACCEPTED; z holds 0 at runtime
}
\end{verbatim}
  is accepted while \texttt{z} holds the number $0$. tsc rejects this
  program; SKJS is unsound here, not merely more permissive in an
  annotation corner.
\item \emph{Spurious rejection.} For \texttt{x: 0.0 | "a"} the union
  appears to have \emph{no} falsy member, and \texttt{if (x)} is rejected
  outright with ``A condition is expected to be of type boolean or a type
  with a single falsy value.'' The same bug, surfacing as incompleteness.
\item \emph{Controls.} The programs above with \texttt{0.0} replaced by
  \texttt{0} behave exactly per the model: \texttt{0 | ""} is rejected as
  a condition (two falsy values); \texttt{0 | "a"} is accepted and
  narrows correctly.
\end{enumerate}

A fourth witness is predicted by code inspection (not yet executed):
\texttt{x: 0 | 0x0} stores a two-element lexeme set whose members are
\emph{both} recognized as zero spellings, so \texttt{falsiness} counts
$\mathsf{FOne} + \mathsf{FOne} = \mathsf{FMany}$ and the condition is
rejected --- although the type is semantically the singleton
$\{0\}$-class and tsc collapses it to the literal type \texttt{0}. This
is the injectivity failure (ii) surfacing independently of the falsiness
failure (iii): here every lexeme is classified \emph{correctly}, and the
duplicate class is still miscounted. Canonicalization fixes it by
collapsing the union at construction.

By contrast the pinned test \texttt{acceptOctalZero}
(\path{tests/projects/truthiness-safe/src/conditions.ts}, lines
106--115) exercises a spelling inside the 14-element stopgap set
(\texttt{0o0 | 1}) and passes today.

%==========================================================================
\section{Denotations $\delta_k$}
\label{sec:denotations}

The main paper assumes a denotation
$\delta_k : \Lit{k} \to \Payload_{\tagfn(k)}$. We now fix it per kind.
$\Lit{k}$ is the set of strings that can be stored for kind $k$ by the
construction sites of Section~\ref{sec:sites}; for num/bigint these are
ECMAScript \emph{NumericLiteral}/\emph{BigIntLiteral} lexemes optionally
prefixed by \texttt{-} (the parser composes the sign), for str the
unquoted token content, for bool $\{\texttt{"true"},\texttt{"false"}\}$,
and for sym synthesized identities.

\begin{definition}[$\delta_{\texttt{num}}$]
\label{def:deltanum}
For a NumericLiteral body $\ell$ (after removing any leading \texttt{-}
and all digit separators \texttt{\_}): its \emph{mathematical value}
$\mathrm{MV}(\ell) \in \mathbb{R}_{\geq 0}$ is given by the ECMA-262
static semantics of NumericLiteral (decimal, or radix $16/8/2$ under
\texttt{0x}/\texttt{0o}/\texttt{0b} prefixes, exponents included), and
$\delta_{\texttt{num}}(\ell)$ is \emph{the Number value for}
$\mathrm{MV}(\ell)$: the IEEE-754 round-to-nearest-ties-to-even image of
that real in $\Fdb$, with overflow to $+\infty$. A leading \texttt{-}
negates the result (giving $-0$ when the body denotes $0$). We also admit
the two words \texttt{Infinity}/\texttt{-Infinity} with denotations
$\pm\infty$ (they arise only as canonical renderings, never as source
lexemes --- \texttt{Infinity} lexes as an identifier).
\end{definition}

\begin{definition}[$\delta_{\texttt{bigint}}$]
For a BigIntLiteral $\ell = m\texttt{n}$ (optionally \texttt{-}-prefixed),
$\delta_{\texttt{bigint}}(\ell) \in \mathbb{Z}$ is the signed integer
value of $m$ in its radix, separators removed. Exact; no rounding.
\end{definition}

\begin{definition}[$\delta_{\texttt{str}}, \delta_{\texttt{bool}},
\delta_{\texttt{sym}}$]
\label{def:deltarest}
$\delta_{\texttt{str}}(\ell)$ is the string value $\ell$ denotes ---
i.e.\ the code-unit sequence \emph{after decoding escape sequences} of
the source string literal. $\delta_{\texttt{bool}}$ maps
\texttt{"true"}/\texttt{"false"} to
$\mathsf{true}/\mathsf{false}$. $\delta_{\texttt{sym}}$ maps each
synthesized position-derived identity
(\texttt{unique\_symbol:file:sym:offset:name}) to a distinct element of
$\SymVal$; as the main paper notes, these are not source lexemes, so
distinct stored strings have distinct denotations \emph{by construction}
--- kind \texttt{sym} is already canonical and plays no further role
here.
\impl{TEnv.sk / uniqueSymbol}
\end{definition}

Two remarks. First, $\delta_{\texttt{str}}$ as defined is the
\emph{semantic} denotation; the implementation currently stores the raw,
\emph{undecoded} content (Finding~\ref{find:escapes}). Second,
$\mathsf{NaN}$ is not in the range of $\delta_{\texttt{num}}$:

\begin{theorem}[No NaN literal]
\label{thm:nonan}
For every $\ell \in \Lit{num}$, $\delta_{\texttt{num}}(\ell) \neq
\mathsf{NaN}$.
\end{theorem}
\begin{proof}
$\mathrm{MV}(\ell)$ is a nonnegative real; the Number value of a real is a
finite double or $\pm\infty$, never $\mathsf{NaN}$, and negation
preserves that. The only non-grammar lexemes admitted are
\texttt{Infinity}/\texttt{-Infinity}, denoting $\pm\infty$. The remaining
mint path, \texttt{keyof} over string keys, filters through
\texttt{isJSNumericString}, whose grammar has no \texttt{NaN} production.
\end{proof}

This \emph{proves} (rather than assumes) the main paper's stipulation
that $\Lit{num}$ contains no $\mathsf{NaN}$ spelling, and shows the
\texttt{literal == "NaN"} arm of \texttt{isFalsy} is dead code today.
\impl{JSNumericString.sk / isJSNumericString; TType.sk / isFalsy}

%==========================================================================
\section{Canonical Forms}
\label{sec:canonical}

The canonicalization strategy is uniform: \emph{parse the lexeme to the
denoted value, then re-render that value by a fixed canonical rendering}.
Injectivity then reduces to properties of the rendering, independent of
how baroque the accepted input grammar is.

\subsection{Numbers}
\label{sec:canon-num}

The canonical rendering is JavaScript's own
$\mathtt{Number::toString}$ (radix 10), written
$\render_{\texttt{num}} : \Fdb \setminus \{\mathsf{NaN}\} \to \String$.
We treat IEEE-754 and the ECMA-262 rendering algorithm as an abstract
interface with three axioms (the ``abstracted-with-axioms'' bin of the
modeling contract); each is a stated property of ECMA-262
\S6.1.6.1.20 (\emph{Number::toString}).

\begin{axiom}[NUM-FUN]
\label{ax:numfun}
$\render_{\texttt{num}}$ is a total function on
$\Fdb \setminus \{\mathsf{NaN}\}$, with
$\render_{\texttt{num}}(\pm\infty) = \texttt{"Infinity"} /
\texttt{"-Infinity"}$.
\end{axiom}

\begin{axiom}[NUM-RT (round-trip)]
\label{ax:numrt}
For every finite $x$,
$\delta_{\texttt{num}}(\render_{\texttt{num}}(x)) \zeq x$. (ECMA-262
chooses digits $s$ and exponent $n$ such that \emph{the Number value for}
$s \cdot 10^{\,n-k}$ \emph{is} $x$; parsing the rendered string recovers
exactly that value. For $x \zeq 0$ the rendering is \texttt{"0"}, which
parses to $+0 \zeq x$.) For $x = \pm\infty$ the round-trip is by the
\texttt{Infinity} clause of Definition~\ref{def:deltanum}.
\end{axiom}

\begin{axiom}[NUM-ZID (zero identification)]
\label{ax:numzid}
$x \zeq y \implies \render_{\texttt{num}}(x) = \render_{\texttt{num}}(y)$.
(The only non-singleton $\zeq$-class is $\{0,-0\}$, and the algorithm
returns \texttt{"0"} for both --- the sign of $-0$ is dropped.)
\end{axiom}

\begin{definition}[Canonical numeric lexemes; $\canon_{\texttt{num}}$]
\label{def:canonnum}
$\Canon{num} \triangleq
\render_{\texttt{num}}(\Fdb \setminus \{\mathsf{NaN}\})$, and
$\canon_{\texttt{num}} \triangleq
\render_{\texttt{num}} \circ\, \delta_{\texttt{num}} : \Lit{num} \to
\Canon{num}$.
\end{definition}

\emph{Justification of the rendering choice.}
(1)~$-0 \mapsto \texttt{"0"}$ (Axiom NUM-ZID) implements the main paper's
$\approx$-closure semantics of literal types --- the literal type
\texttt{0} already denotes $\{0,-0\}$, so a distinct \texttt{-0}
canonical spelling would be a second name for the same type; tsc likewise
gives \texttt{const x = -0} the type \texttt{0}.
(2)~$\mathsf{NaN} \notin \mathrm{dom}(\render_{\texttt{num}})$ matches
the main paper's exclusion of a NaN literal type (its denotation would be
$\emptyset$ under closure semantics), and Theorem~\ref{thm:nonan} shows
totality is not thereby lost.
(3)~The finite part of $\Canon{num}$ is exactly the strings satisfying
$\mathtt{String(Number(s))} = s$, the predicate the codebase already
uses to classify numeric property names --- so \texttt{keyof}-minted
number literals and canonicalized literals agree by construction (but
see Finding~\ref{find:jsnum} for a defect in the current implementation
of that predicate).
(4)~It is what tsc renders in hovers and errors, minimizing user-visible
divergence.
\impl{JSNumericString.sk / isJSNumericString (intended grammar)}

Note $\Canon{num}$ contains renderings such as \texttt{"1e-7"} whose
signed exponent the current SKJS lexer cannot tokenize as a single
numeric token (Section~\ref{sec:validation}); this is harmless because
stored lexemes are never re-lexed, and tsc accepts these spellings.

\subsection{Bigints}
\label{sec:canon-bigint}

Here everything is concrete; no axioms are needed.

\begin{definition}[Canonical bigint lexemes]
\label{def:canonbigint}
$\render_{\mathbb{Z}}(z)$ is the standard signed decimal numeral of
$z \in \mathbb{Z}$ (no leading zeros, no separators, no radix prefix, no
sign for $z \geq 0$, so $\render_{\mathbb{Z}}(0) = \texttt{"0"}$ ---
in particular \texttt{-0n} normalizes to \texttt{0n}), and
$\Canon{bigint} \triangleq \{\render_{\mathbb{Z}}(z)\,\texttt{n} \mid z
\in \mathbb{Z}\}$, with
$\canon_{\texttt{bigint}} = (\render_{\mathbb{Z}} \circ\,
\delta_{\texttt{bigint}})\,\texttt{n}$.
\end{definition}

\begin{fact}
\label{fact:decimal}
$\render_{\mathbb{Z}}$ is a bijection $\mathbb{Z} \to
\render_{\mathbb{Z}}(\mathbb{Z})$, and
$\delta_{\texttt{bigint}}(\render_{\mathbb{Z}}(z)\,\texttt{n}) = z$.
(Uniqueness of positional decimal numerals without leading zeros;
elementary.)
\end{fact}

\subsection{Strings, booleans, symbols}
\label{sec:canon-rest}

$\Canon{str} \triangleq \String$: the canonical form of a string literal
is \emph{the denoted string itself} ($\canon_{\texttt{str}} =
\delta_{\texttt{str}}$, i.e.\ decode escapes at the quote boundary and
store code units). $\Canon{bool} = \{\texttt{"true"},
\texttt{"false"}\}$ and $\Canon{sym} = $ the synthesized identities:
already canonical, $\canon$ is the identity on them
(Definition~\ref{def:deltarest}).

\begin{definition}[Canonical-storage invariant]
\label{def:invariant}
A \texttt{TPrimValue(kind $k$, Some(lits))} is \emph{canonical} when
every stored lexeme is in $\Canon{k}$. The fix of Section~\ref{sec:fix}
establishes this as a representation invariant of the type layer.
\end{definition}

%==========================================================================
\section{Proofs over Canonical Forms}
\label{sec:proofs}

Throughout this section literals are canonical
(Definition~\ref{def:invariant}); kind \texttt{sym} is covered by
construction and omitted from the case analyses where it is trivial.

\begin{lemma}[Idempotence]
\label{lem:idem}
For $k \in \{\texttt{num}, \texttt{bigint}\}$ and $\ell \in \Canon{k}$:
$\canon_k(\ell) = \ell$.
\end{lemma}
\begin{proof}
\texttt{num}: $\ell = \render(x)$ for some $x$; then
$\canon(\ell) = \render(\delta(\render(x)))$, and by NUM-RT
$\delta(\render(x)) \zeq x$, so by NUM-ZID
$\render(\delta(\render(x))) = \render(x) = \ell$ (infinite $x$: direct).
\texttt{bigint}: immediate from Fact~\ref{fact:decimal}.
\end{proof}

\begin{theorem}[L-LIT(i): $\delta$ is a function; totality]
\label{thm:fun}
Each $\delta_k$ is a well-defined function on $\Lit{k}$, and each
$\canon_k$ is a well-defined function $\Lit{k} \to \Canon{k}$. For
\texttt{str} this holds with $\delta_{\texttt{str}}$ as specified in
Definition~\ref{def:deltarest}; the implementation currently realizes a
different function (raw storage), see Finding~\ref{find:escapes}.
\end{theorem}
\begin{proof}
\texttt{num}: MV is defined by a deterministic grammar attribution;
rounding is a function; Axiom NUM-FUN gives the rendering
(Theorem~\ref{thm:nonan} keeps $\mathsf{NaN}$ out of its domain).
\texttt{bigint}: radix valuation and Fact~\ref{fact:decimal}.
\texttt{str}: escape decoding is deterministic. \texttt{bool},
\texttt{sym}: finite/by-construction tables.
\end{proof}

\begin{theorem}[L-LIT(ii): injectivity modulo $\approx$]
\label{thm:inj}
For every kind $k$ and $\ell_1, \ell_2 \in \Canon{k}$: if
$(\tagfn(k), \delta_k(\ell_1)) \approx (\tagfn(k), \delta_k(\ell_2))$
then $\ell_1 = \ell_2$. Equivalently: distinct canonical lexemes denote
$\approx$-inequivalent values, so each $\approx$-class of payloads meets
at most one canonical lexeme.
\end{theorem}
\begin{proof}
\texttt{num}: on num payloads $\approx$ is $\zeq$. Write
$\ell_i = \render(x_i)$. By NUM-RT $\delta(\ell_i) \zeq x_i$
(for infinite $x_i$, $\delta(\ell_i) = x_i$), so
$\delta(\ell_1) \zeq \delta(\ell_2)$ gives $x_1 \zeq x_2$ by
symmetry/transitivity of $\zeq$ on its domain ($\mathsf{NaN}$ is not in
the range of $\delta$ by Theorem~\ref{thm:nonan}). NUM-ZID then gives
$\ell_1 = \render(x_1) = \render(x_2) = \ell_2$.
\texttt{bigint}: $\approx$ is equality on $\mathbb{Z}$;
Fact~\ref{fact:decimal} (bijectivity) concludes.
\texttt{str}: $\canon_{\texttt{str}}$ stores the value itself, so
$\delta$ restricted to $\Canon{str}$ is the identity --- injective.
\texttt{bool}: two-point table, injective. \texttt{sym}: by construction
(main paper, Definition ``Literal Denotation'').
\end{proof}

\begin{theorem}[L-LIT(iii): falsiness-exactness]
\label{thm:falsy}
Define the \emph{falsy canonical sets}
\[
F_{\texttt{bool}} = \{\texttt{"false"}\},\quad
F_{\texttt{num}} = \{\texttt{"0"}\},\quad
F_{\texttt{str}} = \{\texttt{""}\},\quad
F_{\texttt{bigint}} = \{\texttt{"0n"}\},\quad
F_{\texttt{sym}} = \emptyset .
\]
For every kind $k$ and $\ell \in \Canon{k}$:
$\truthy(\tagfn(k), \delta_k(\ell)) = \mathsf{false}
\iff \ell \in F_k$. Hence over canonical forms the lexeme test
``$\ell \in F_k$'' agrees exactly with the main paper's Definition of
Truthiness, and each $F_k$ is finite (at most one element).
\end{theorem}
\begin{proof}
By the Definition of Truthiness, the falsy payloads are:
$\mathsf{false}$ (bool); $\{0, -0, \mathsf{NaN}\}$ (num); \texttt{""}
(str); $0$ (bigint); none (sym).
\texttt{num}: $\mathsf{NaN} \notin \mathrm{range}(\delta)$
(Theorem~\ref{thm:nonan}); $\delta(\ell) \in \{0,-0\}$ iff
$\delta(\ell) \zeq 0$ iff (by Lemma~\ref{lem:idem} and NUM-ZID)
$\ell = \canon(\ell) = \render(\delta(\ell)) = \render(0) =
\texttt{"0"}$; conversely $\delta(\texttt{"0"}) = 0$ is falsy.
\texttt{bigint}: $\delta(\ell) = 0$ iff $\ell = \texttt{"0n"}$ by
Fact~\ref{fact:decimal}. \texttt{str}: identity denotation, falsy iff the
stored string is empty. \texttt{bool}, \texttt{sym}: by the tables.
\end{proof}

\begin{corollary}[Residue lemma discharged]
\label{cor:residue}
With canonical storage and \texttt{isFalsy} replaced by membership in
$F_k$ (a single string equality per kind), the literal-filtering case of
the main paper's Lemma ``Residues are Sound'' holds unconditionally: the
kept-literal sets of $\mathsf{tres}_b$ and $\mathsf{vres}_b$ contain
exactly the literals whose denotations have truthiness $b$.
\impl{TType.sk / truthinessResidue, truthinessValueResidue}
\end{corollary}

\begin{corollary}[Exactness of constant removal]
\label{cor:rmc}
The exactness clause of the main paper's Theorem ``Soundness of Constant
Removal'' holds under canonical storage: its stated side condition is
precisely L-LIT(ii), Theorem~\ref{thm:inj}. (Its soundness clauses R1/R2
were unconditional already.)
\impl{TType.sk / removeConstant; PrimLiterals.sk / diff}
\end{corollary}

\begin{corollary}[Lexeme set operations compute value set operations]
\label{cor:setops}
For canonical $L_1, L_2 \subseteq \Canon{k}$, writing
$\sem{L} \triangleq \sem{\mathsf{prim}(k, L)}$:
$\sem{L_1 \cup L_2} = \sem{L_1} \cup \sem{L_2}$ (unconditional);
$\sem{L_1 \cap L_2} = \sem{L_1} \cap \sem{L_2}$; and
$\sem{L_1 \setminus L_2} = \sem{L_1} \setminus \sem{L_2}$.
\end{corollary}
\begin{proof}
Union: $\sem{\cdot}$ is a union of $\approx$-closures, which commutes
with $\cup$. Intersection and difference: by Theorem~\ref{thm:inj} the
$\approx$-classes named by distinct canonical lexemes are disjoint, so a
value $v \approx \delta(\ell_1) $ with $\ell_1 \in L_1$ lies in
$\sem{L_2}$ iff $\ell_1 \in L_2$; the class-wise identities follow.
(Without canonicity only $\subseteq$ / $\supseteq$ over-approximations
hold, which is why \texttt{PrimLiterals::inter} can spuriously produce
$\bot$ today: $\{\texttt{"0"}\} \cap \{\texttt{"0x0"}\} = \emptyset$ on
lexemes while the value sets coincide --- an unsound \emph{emptiness}
verdict wherever $\bot$ is read as disjointness.)
\impl{PrimLiterals.sk / union, inter, diff}
\end{proof}

\begin{corollary}[Exact falsy counting]
\label{cor:count}
Over canonical storage, \texttt{SType::falsiness} counts, per literal
member, the exact number of falsy $\approx$-classes ($0$ or $1$, since
$|F_k| \leq 1$), so \texttt{FalsyCount} never double-counts a class
across spellings (the \texttt{0 | 0x0} witness of
Section~\ref{sec:witnesses} disappears) and the \texttt{safe}-mode
condition guard implements its intended semantics on literal unions.
\impl{TType.sk / SType::falsiness, FalsyCount;
TypeCheckUtils.sk / checkCondition}
\end{corollary}

%==========================================================================
\section{Fix Specification}
\label{sec:fix}

\subsection{Where normalization lives}
\label{sec:where}

Two placements were considered.

\paragraph{(a) In the parser (\texttt{getNumber}/\texttt{getBigInt}).}
One choke point, and it would discharge the two \texttt{TODO: check and
normalize somehow} comments in place. Rejected as the \emph{sole}
placement for two reasons: (1)~those helpers also feed the AST consumed
by non-type clients --- numeric \emph{property names}
(\texttt{ParserDefs.sk / getNameOrStringOrNumber}) and diagnostic text
quoting source --- so normalizing there silently changes behaviors that
deserve their own decision (Section~\ref{sec:propnames}); (2)~the
negative-literal sites compose \texttt{"-"} with the lexeme \emph{after}
these helpers return (\texttt{ParserType.sk} lines 232--233;
\texttt{directLiteralSType}; \texttt{asConstExpr}), so parser-side
normalization alone cannot maintain the invariant --- \texttt{-0} would
still reach the type layer as \texttt{"-0"}.

\paragraph{(b) In the type layer, at literal-type construction
(adopted).}
Introduce
\begin{center}
\texttt{fun canonLexeme(kind: TPrimValueKind, raw: String): Result<String,
CanonError>}
\end{center}
implementing $\canon_k$ (num: strip separators, parse sign/radix/MV,
round, render per \S\ref{sec:canon-num}; bigint: per
\S\ref{sec:canon-bigint}; str/bool/sym: identity), and route \emph{every}
constructor of a num/bigint \texttt{PrimLiteral} through it:
\begin{itemize}
\item \texttt{TPrimValue::single} (TType.sk ${\sim}$1521): normalize when
  \texttt{kind} is \texttt{TNum()}/\texttt{TBigInt()}. This covers
  \texttt{TEnv.sk} 1054/1205 (annotations, including parser-composed
  negatives), \texttt{directLiteralSType}, and \texttt{FieldKey::toType}.
\item \texttt{asConstExpr}'s \texttt{UOMinus} arm
  (TypeCheckExpr.sk ${\sim}$2803): it sign-flips by string prefix; on
  canonical input \texttt{"-"}${}+{}$canonical is canonical
  \emph{except} for \texttt{"0"}/\texttt{"0n"} (yielding
  \texttt{"-0"}/\texttt{"-0n"}), so the arm must re-normalize --- funneling
  through \texttt{TPrimValue::single} suffices.
\item \texttt{PrimLiterals::ofSet} call sites with numeric kind: the
  \texttt{keyof} site (\texttt{TEnv.sk} 2162) already receives canonical
  strings \emph{by intent} (they pass \texttt{isJSNumericString}), but
  Finding~\ref{find:jsnum} shows that predicate over-accepts; route
  through \texttt{canonLexeme} (or fix the predicate) rather than assume.
\item Already-canonical constants need no change, only the documented
  invariant: \texttt{SyntheticTuples.sk} 127 (\texttt{Int.toString}),
  \path{truthinessValueResidue}'s \texttt{"0n"}/\texttt{""}, boolean
  literals.
\end{itemize}
Document Definition~\ref{def:invariant} on \texttt{PrimLiterals} (``for
num/bigint kinds, \texttt{val} is canonical''), ideally with a debug
assertion \texttt{canonLexeme(kind, val) == val} at construction.
\impl{TType.sk / TPrimValue::single; PrimLiterals.sk / PrimLiterals::ofSet}

\subsection{Validation of over-lexed input}
\label{sec:validation}

The numeric-token loop in the lexer accepts a strict superset of the
NumericLiteral grammar --- any run over
$\{\texttt{0-9}, \texttt{a-f}, \texttt{A-F}, \texttt{e}, \texttt{E},
\texttt{.}, \texttt{x}, \texttt{X}, \texttt{b}, \texttt{B}, \texttt{o},
\texttt{O}, \texttt{\_}\}$ --- so lexemes like \texttt{1.2.3},
\texttt{0x}, or \texttt{1\_\_2} reach \texttt{getNumber} and are stored
raw today. \texttt{canonLexeme} is therefore \emph{partial} on raw input
and must return a diagnostic on non-grammar lexemes; this discharges the
``check'' half of the parser TODOs. Separately, signed exponents
(\texttt{1e-5}) never arrive as one token (the loop stops at
\texttt{-}/\texttt{+}), so the normalizer's input grammar has unsigned
exponents only, while its \emph{output} may contain signed ones
(\texttt{"1e-7"}); recommend an independent lexer issue for signed
exponents.
\impl{Lexer.sk / numeric-token loop (${\sim}$384--412);
ParserUtils.sk / getNumber, getBigInt}

\subsection{User-visible pretty-printing}

Literal types print their stored lexemes
(\texttt{PrimLiterals::toStringList}), so diagnostics and hovers change
spelling: \texttt{0o0 | 1} prints as \texttt{0 | 1}, \texttt{0x1 | 1}
collapses to \texttt{1}. This matches tsc's rendering. Two corners:
(1)~large/small magnitudes print in canonical scientific form
(\texttt{1e+21}, \texttt{1e-7}) --- valid TS literal spellings, though
SKJS's own lexer cannot yet re-lex the signed-exponent ones
(\S\ref{sec:validation}); (2)~overflowed literals print as
\texttt{Infinity}, which is \emph{not} a literal-type spelling --- either
render as a spelling that denotes it (\texttt{1e999}) or accept the
tsc-divergent print; we recommend the former for round-trippable output.
\impl{PrimLiterals.sk / toStringList; TTypePp.sk line 210}

\subsection{Interaction with \texttt{PrimLiterals} operations}

No code change: \texttt{union}/\texttt{inter}/\texttt{diff}/%
\texttt{contains}/\texttt{xxhash} remain lexeme-level; the invariant
upgrades their \emph{meaning} to exact value-set operations
(Corollary~\ref{cor:setops}). In particular \texttt{sameConstant} and
structural type hashing become value-respecting, and
\texttt{TFields::isTuple}'s \texttt{toIntOption} on the \texttt{length}
literal starts working for previously radix-spelled lengths
(\texttt{"0x2"} fails \texttt{toIntOption} today; canonical \texttt{"2"}
succeeds) --- a strict precision improvement.
\impl{PrimLiterals.sk; TType.sk / TPrim::sameConstant,
TFields::isTuple (${\sim}$2105)}

\subsection{Numeric property names (companion change)}
\label{sec:propnames}

Numeric \emph{property names} take the same raw-lexeme path
(\path{ParserDefs.sk / getNameOrStringOrNumber}, lines 64--71) and end
up as field-name strings compared against literal-type lexemes
(\path{FieldKey::ofSType}; \path{TFields::maybeGet};
\path{namedFieldsForKeyKind} filters names by
\path{isJSNumberKeyName}). TS defines a numeric property name as
$\mathtt{String(Number(name))}$, i.e.\ exactly $\Canon{num}$'s finite
part. For coherence, a declaration \texttt{\{ 0x10: T \}} must store the
field name \texttt{"16"} so that the canonicalized access key
\texttt{obj[0x10]} (literal type \texttt{16}) resolves to it ---
otherwise the fix would leave declaration-side and use-side spellings
permanently disconnected. This is the one place parser-side
normalization \emph{is} required; it is a separable, behavior-visible
companion change and should ship with the same issue.
\impl{ParserDefs.sk / getNameOrStringOrNumber; TType.sk /
FieldKey::ofSType, namedFieldsForKeyKind}

\subsection{What becomes deletable}

\begin{itemize}
\item \texttt{TPrimValueKind::isZeroSpelling} (TType.sk ${\sim}$1242):
  deleted entirely, with the \#655 stopgap comment (${\sim}$1235).
\item \texttt{isFalsy} (${\sim}$1224) collapses to the falsy canonical
  sets of Theorem~\ref{thm:falsy}: \texttt{TNum() -> literal == "0"}
  (the \texttt{"NaN"} arm is dead by Theorem~\ref{thm:nonan});
  \texttt{TBigInt() -> literal == "0n"} (no \texttt{stripSuffix});
  other kinds unchanged.
\item The two parser \texttt{TODO}s in \texttt{ParserUtils.sk} are
  resolved by \S\ref{sec:validation}.
\end{itemize}

\subsection{Migration notes}

A repository-wide search at \texttt{9108f339} found exactly one test
depending on a non-canonical spelling: \texttt{acceptOctalZero}
(\texttt{tests/projects/truthiness-safe/src/conditions.ts} 106--115),
which annotates with \texttt{0o0}. It \emph{survives} the fix --- both
the union member and the else-branch annotation \texttt{const e: 0o0}
normalize to \texttt{0}, so the assignment still checks --- but its
comment describes the stopgap and should be rewritten to describe
canonicalization. The test suites are \texttt{@skjs-expect-error}
annotation-based with no expected-output snapshots, so no golden files
capture literal-type rendering; error-message \emph{text} changes
(canonical spellings) affect only tests that grep message contents, of
which none were found for radix/float zeros.

%==========================================================================
\section{Findings from Failed and Partial Proof Attempts}
\label{sec:findings}

Per the series contract, unprovable steps are reported as findings, not
silently narrowed away.

\begin{finding}[String escapes: L-LIT(ii) fails for kind \texttt{str},
witnessed]
\label{find:escapes}
Attempting to prove Theorem~\ref{thm:inj} for \texttt{str} \emph{as
implemented} fails. The lexer skips escape sequences without decoding
them (\texttt{Lexer.sk} string loops, ${\sim}$239--277: on
\texttt{\textbackslash} it advances two positions), and
\texttt{unquoteString} only strips the quotes, with an explicit
\texttt{TODO: handle escaped sequences}. So the literal types written
\texttt{"\textbackslash x61"} and \texttt{"a"} store \emph{distinct}
lexemes denoting the \emph{same} runtime string --- a second violation of
L-LIT(ii), independent of \#655 and not covered by its fix. Consequences
mirror the numeric case: \texttt{inter} can return $\bot$ for
overlapping string literal types (unsound where $\bot$ is read as
disjointness), \texttt{diff}/exactness degrade, and equality narrowing
against an escape-spelled constant misclassifies. Falsiness is
\emph{also} affected, contrary to an earlier belief that no non-empty
lexeme can denote \texttt{""}: ECMA-262 \emph{LineContinuation}
(backslash + line terminator) has string value \texttt{""}, so the
non-empty backslash-newline lexeme denotes the
falsy \texttt{""} while \texttt{isFalsy} (\texttt{lexeme == ""}) reads
it truthy. This is a string-kind violation of \textbf{L-LIT}(iii),
witnessed end-to-end: a union \texttt{"\ldots" | 0} whose string member
is a line-continuation lexeme is accepted with \texttt{const z:\ 0 = x}
in the \texttt{else} of \texttt{if (x)} while that branch runs at
\texttt{x = ""} (the canonical \texttt{"" | 0} control is correctly
rejected). The injectivity failure
is witnessed too: with
\texttt{x: "a" | 0}, the test
\texttt{x === "\textbackslash x61"} types its (runtime-taken) true branch
\texttt{never}, and \texttt{const z: 0 = x} is accepted while \texttt{z}
holds the string \texttt{"a"}; the checker even emits an always-false
warning for a comparison that is true at runtime. Issue draft:
\path{issues/string-literal-escapes-undecoded.md}. Until decoding
lands, Theorem~\ref{thm:inj}'s \texttt{str} case holds only under the
assumption ``stored lexeme $=$ denoted string,'' which is exactly
$\canon_{\texttt{str}}$ of \S\ref{sec:canon-rest}.
\impl{ParserUtils.sk / unquoteString; Lexer.sk / string loops}
\end{finding}

\begin{finding}[\texttt{isJSNumericString} over-accepts at the exponent
thresholds]
\label{find:jsnum}
We attempted to prove that the finite part of $\Canon{num}$
(Definition~\ref{def:canonnum}) equals the grammar implemented by
\texttt{isJSNumericString} --- its header claims exactly the round-trip
characterization $\mathtt{String(Number(s))} = s$. The proof fails: the
implemented scientific production
$[\texttt{-}][\texttt{1-9}](.\,d^{*})\texttt{e}[\pm]e^{+}$
omits ECMA-262's magnitude side condition (scientific rendering is used
only when the printed exponent is $\geq 21$ or $\leq -7$).
Witness: \texttt{"1e+5"} is accepted, but
$\mathtt{String(Number("1e{+}5"))} = \texttt{"100000"} \neq
\texttt{"1e+5"}$, so it is not canonical. Impact today: (a)~the
\texttt{keyof} path can mint the non-canonical number literal
\texttt{"1e+5"} from a string key, and (b)~number-index governance
treats such field names as reachable by number keys though no
\texttt{String(n)} produces them. Impact on the fix: if
\texttt{isJSNumericString} were used as the canonical-form
\emph{validator}, injectivity (Theorem~\ref{thm:inj}) would fail on
\texttt{"1e+5"} vs.\ \texttt{"100000"}. The fix must therefore either
add the threshold conditions to the grammar or route these strings
through \texttt{canonLexeme}; we recommend both (the predicate has
independent index-signature clients). $\Canon{num}$ in this paper is
defined by the rendering, not by the predicate, precisely to keep the
theorems immune to this defect. Adversarial probing of impacts (a) and
(b) found \emph{no soundness leak} — every exploitation attempt (index
reads and writes through \texttt{"1e+5"}, returning or indexing with the
minted \texttt{keyof} member) is rejected by the checker. The defect is
real but its observable effects are completeness/classification only; it
is a hygiene fix, not an unsoundness, and accordingly has no issue draft
in \texttt{issues/}.
\impl{JSNumericString.sk / isJSNumericString,
isJSNumericExponentPartIter; TEnv.sk line 2162; TType.sk line 1794}
\end{finding}

\begin{finding}[Rejected alternative: compare-by-value at use sites]
\label{find:usesite}
An alternative fix --- keep raw lexemes and parse to values inside
\texttt{isFalsy}, \texttt{inter}, \texttt{diff}, \texttt{sameConstant},
\texttt{xxhash}, subtyping, and printing --- was considered and
rejected. The comparison surface is wide and grows (any future consumer
of \texttt{PrimLiteral.val} silently reverts to lexeme semantics), the
hash requires a canonical key anyway (hashing raw spellings breaks
hash-consistency of semantically equal types), and per-use parsing
re-does work the constructor does once. Normalize-at-construction makes
the invariant local and every downstream operation correct by
Corollary~\ref{cor:setops} without modification. This mirrors the issue
tracker's own recommendation on \#655.
\end{finding}

%==========================================================================
\section{Status and Ledger Edits}
\label{sec:status}

\subsection{Effect of the fix}

\begingroup
\renewcommand{\arraystretch}{1.25}
\small
\begin{longtable}{@{}p{2.4cm}p{11.0cm}@{}}
\toprule
\textbf{Item} & \textbf{After the fix of \S\ref{sec:fix}} \\
\midrule
\endhead
L-LIT(i) & unchanged (already established: $\delta_k$ concrete per kind,
  \S\ref{sec:denotations}; NaN-freeness of $\Lit{num}$ a theorem,
  Thm.~\ref{thm:nonan}); the \texttt{str} raw-storage caveat
  (Finding~\ref{find:escapes}) persists until escape decoding \\
L-LIT(ii) & closes for num/bigint/bool/sym; \texttt{str} conditional on
  escape decoding (injectivity over canonical forms, Thm.~\ref{thm:inj};
  the string-kind violation is witnessed end-to-end,
  Finding~\ref{find:escapes}, issue draft
  \path{issues/string-literal-escapes-undecoded.md}) \\
L-LIT(iii) & closes for num/bigint/bool/sym; for \texttt{str},
  conditional on escape \emph{and} line-continuation decoding
  (Finding~\ref{find:escapes} witnesses both an (ii) and an (iii)
  string-kind violation) \\
Residue lemma (main paper) & unconditional (Cor.~\ref{cor:residue}) \\
rmc exactness (main paper) & unconditional (Cor.~\ref{cor:rmc}) \\
New obligations & \texttt{canonLexeme} correctness
  ($= \canon_k$ of \S\ref{sec:canonical}); property-name
  normalization (\S\ref{sec:propnames}); \texttt{isJSNumericString}
  thresholds (Finding~\ref{find:jsnum}) --- each small and pinnable \\
\bottomrule
\end{longtable}
\endgroup

Until the fix lands, the main paper's \textsc{violated} marking on
L-LIT(ii)/(iii) stands unchanged; nothing in this paper weakens it.

\subsection{Registry}

The authoritative, current statuses are maintained in the registry
(\texttt{docs/ledger.yaml}), rendered as the main paper's Assumptions and
Proof Obligations tables; this paper is cited there as evidence. The
registry rows it supports:

% GENERATED by docs/tools/docs.py from docs/ledger.yaml — do not edit.

\begin{center}
\renewcommand{\arraystretch}{1.25}
\small
\begin{tabular}{@{}>{\raggedright\arraybackslash}p{2.4cm}>{\raggedright\arraybackslash}p{8.2cm}>{\raggedright\arraybackslash}p{2.6cm}@{}}
\toprule
\textbf{Row} & \textbf{Statement} & \textbf{Status} \\
\midrule
L-LIT & Literal spellings: (i) $\delta$ is a function; (ii) injectivity modulo $\approx$ after normalization; (iii) falsiness-exactness of the lexeme test. & \textsc{violated} \\
\bottomrule
\end{tabular}
\end{center}


\paragraph{Imports (review contract).}
The named statements this paper assumes from elsewhere. A reviewer treats
each as an axiom, checking citations against the registry's canonical
statement; no other paper's internals are in scope:

% GENERATED by docs/tools/docs.py from docs/ledger.yaml — do not edit.

\noindent\emph{None beyond the main paper's definitions.}


\subsection{Registry impact of the fix}

When the canonicalization fix of Section~\ref{sec:fix} lands, the
registry changes as follows.
\begin{enumerate}
\item \textbf{L-LIT status.} (i) \textsc{proved}; (ii)/(iii)
  \textsc{proved} (this paper, Thms.~\ref{thm:inj}/\ref{thm:falsy})
  for num/bigint/bool/sym under the canonical-storage invariant, pinned
  by the tests below; for \texttt{str}, (ii) \textsc{assumed} pending
  escape decoding (its own issue). Anchor field: \texttt{isZeroSpelling}
  (deleted by the fix) is replaced by \texttt{canonLexeme}.
\item \textbf{Main paper, worked example ``Known unsoundness
  (\#655)''.} Becomes a discharged example whose pins flip polarity
  (see below); the generated summary's \textsc{violated} grouping drops
  its L-LIT items (except the string-kind escape entry, until decoding
  lands).
\item \textbf{Main paper, Definition ``Literal Denotation''.} The
  assumption ``$\Lit{num}$ contains no NaN spelling'' becomes a citation
  of Theorem~\ref{thm:nonan}.
\end{enumerate}

\subsection{Recommended pinned tests}

To ship with the fix, in
\texttt{tests/projects/truthiness-safe/src/conditions.ts} unless noted:
\begin{enumerate}
\item \texttt{x: 0.0 | ""} --- \texttt{if (x)} becomes
  \texttt{@skjs-expect-error} (two falsy values), matching today's
  \texttt{0 | ""} control; this inverts the current unsound acceptance
  of \texttt{const z: "" = x} in the else-branch
  (witness~1, \S\ref{sec:witnesses}).
\item \texttt{x: 0.0 | "a"} --- \texttt{if (x)} accepted; then-branch
  narrows to \texttt{"a"}, else-branch to \texttt{0} (and
  \texttt{const e: 0.0 = x} checks, both spellings canonicalizing);
  inverts witness~2.
\item \texttt{x: 0 | 0x0} --- accepted as a condition (the union
  collapses to one falsy class); pins Corollary~\ref{cor:count} and the
  predicted fourth witness.
\item Keep \texttt{acceptOctalZero} with an updated comment
  (\S migration).
\item A rendering pin: a diagnostic that prints a literal type declared
  \texttt{0o0 | 1} must show \texttt{0 | 1}.
\item (With Finding~\ref{find:jsnum}) \texttt{keyof} over
  \texttt{\{ "1e+5": T \}} must \emph{not} yield a numeric literal
  member.
\item (With Finding~\ref{find:escapes}, separate issue) an equality or
  assignment mixing \texttt{"a"} and \texttt{"\textbackslash x61"}
  literal types, pinned to current (incorrect) behavior until decoding
  lands.
\end{enumerate}

Anchor cross-check hooks per the main paper's maintenance appendix: the
deleted \path{isZeroSpelling} anchor must be removed from the L-LIT
row in the same change that deletes the function, and
\texttt{canonLexeme} becomes the row's anchor.

\end{document}
